\section{El problema de las actualizaciones de gradiente en varios pasos}

On-policy actor-critic solo puede hacer un paso de actualización de gradiente por cada lote de datos recolectado, lo cual es muy poco eficiente. Queremos hacer \textbf{varios pasos} de actualización sobre el mismo lote de datos.

\subsection{Repaso de los Importance Weights}

Al usar datos de la política antigua $\pi_\theta$ para actualizar la política nueva $\pi_{\theta'}$, se necesitan importance weights:
$$\nabla_{\theta'} J(\theta') \approx \sum_{t,i} \frac{\pi_{\theta'}(a_{i,t} \mid s_{i,t})}{\pi_\theta(a_{i,t} \mid s_{i,t})} \nabla_{\theta'} \log \pi_{\theta'}(a_{t,i} \mid s_{t,i}) \hat{A}^{\pi_\theta}(s_{t,i}, a_{t,i})$$

\begin{warningbox}{Riesgo de las actualizaciones en varios pasos}
Si se hacen muchos pasos de ascenso de gradiente sobre la función objetivo sustituta (surrogate objective), la política se ve incentivada a alejarse mucho de la política antigua. En ese caso los importance weights dejan de ser confiables, la estimación de la ventaja también deja de ser válida, y esto provoca sobreajuste y colapso del desempeño.
\end{warningbox}

\subsection{Resumen de la sección}

Las actualizaciones de gradiente en varios pasos mejoran la eficiencia de los datos, pero es necesario restringir la magnitud de la actualización de la política para mantener la estabilidad.

